
\documentclass[12pt,letterpaper]{article}

% ---------------- PAQUETES ----------------
\usepackage[spanish,es-tabla]{babel}
\usepackage[utf8]{inputenc}
\usepackage[T1]{fontenc}

\usepackage[left=2.5cm,right=2.5cm,top=2.5cm,bottom=2.5cm]{geometry}
\usepackage{graphicx}
\usepackage{fancyhdr}
\usepackage{titlesec}
\usepackage{xcolor}
\usepackage{booktabs}
\usepackage{longtable}
\usepackage{array}
\usepackage{parskip}
\usepackage{enumitem}
\usepackage{amssymb}
\usepackage{hyperref}

\begin{document}

\section{Guion Técnico de Apoyo para Exposición (Audiencia Experta)}

\subsection{Arquitectura General}

\begin{itemize}
    \item Arquitectura desacoplada tipo cliente--API.
    \item Separación estricta entre:
    \begin{itemize}
        \item Capa de presentación (HTML/CSS/JS).
        \item Capa de servicios (API REST).
    \end{itemize}
    \item API implementada en FastAPI:
    \begin{itemize}
        \item Tipado fuerte.
        \item Serialización automática.
        \item OpenAPI nativo.
    \end{itemize}
    \item Persistencia en PostgreSQL mediante SQLAlchemy ORM.
\end{itemize}

\subsection{Capa de Datos}

\subsubsection{Configuración (\texttt{database.py})}

\begin{itemize}
    \item Engine con pool de conexiones.
    \item \texttt{SessionLocal} como unidad transaccional.
    \item Inyección de sesión en endpoints.
\end{itemize}

\subsubsection{Modelo (\texttt{models.py})}

\begin{itemize}
    \item Entidad \texttt{User} mapeada con SQLAlchemy.
    \item Restricción de unicidad:
    \begin{itemize}
        \item \texttt{correo} con \texttt{unique=True}.
    \end{itemize}
    \item Contraseña:
    \begin{itemize}
        \item Nunca se persiste en claro.
        \item Solo hash persistente.
    \end{itemize}
\end{itemize}

\subsection{Capa de Esquemas}

\subsubsection{Validación (\texttt{schemas.py})}

\begin{itemize}
    \item Uso de Pydantic como frontera semántica.
    \item DTOs:
    \begin{itemize}
        \item \texttt{UserCreate}.
        \item \texttt{UserUpdate}.
    \end{itemize}
    \item Validación previa a:
    \begin{itemize}
        \item Lógica de negocio.
        \item Acceso a datos.
    \end{itemize}
\end{itemize}

\subsection{Capa de Seguridad}

\subsubsection{Criptografía (\texttt{security.py})}

\begin{itemize}
    \item Hashing con bcrypt (passlib).
    \item Tokens JWT:
    \begin{itemize}
        \item Payload: identidad + rol.
        \item Expiración definida.
    \end{itemize}
    \item Autorización vía dependencia:
    \begin{itemize}
        \item \texttt{verify\_token}.
    \end{itemize}
\end{itemize}

\subsection{Lógica de Negocio}

\subsubsection{CRUD (\texttt{crud.py})}

\begin{itemize}
    \item Acceso directo a ORM.
    \item Función crítica:
    \begin{itemize}
        \item \texttt{actualizar\_usuario}.
    \end{itemize}
    \item Uso de:
    \begin{itemize}
        \item \texttt{exclude\_unset=True}.
    \end{itemize}
    \item Implica:
    \begin{itemize}
        \item Actualización parcial.
        \item No sobreescritura involuntaria de credenciales.
    \end{itemize}
\end{itemize}

\subsubsection{Controlador (\texttt{main.py})}

\begin{itemize}
    \item Orquestación de dependencias.
    \item Implementación de RBAC.
    \item Roles privilegiados:
    \begin{itemize}
        \item director.
        \item coordinador.
    \end{itemize}
    \item Restricción mediante:
    \begin{itemize}
        \item Evaluación del rol desde JWT.
        \item Excepción HTTP 403.
    \end{itemize}
\end{itemize}

\subsection{Frontend}

\subsubsection{Estilo (\texttt{style.css})}

\begin{itemize}
    \item Neomorfismo:
    \begin{itemize}
        \item Doble \texttt{box-shadow}.
        \item Ilusión de relieve.
    \end{itemize}
\end{itemize}

\subsubsection{Cliente (\texttt{login.js}, \texttt{dashboard.js})}

\begin{itemize}
    \item Consumo REST vía Fetch API.
    \item Persistencia de sesión:
    \begin{itemize}
        \item Token en \texttt{localStorage}.
    \end{itemize}
    \item Filtro local:
    \begin{itemize}
        \item \texttt{filtrarTabla()}.
        \item Evita round-trips innecesarios.
    \end{itemize}
    \item Control visual de permisos:
    \begin{itemize}
        \item UI reactiva al rol.
        \item Seguridad real delegada al backend.
    \end{itemize}
\end{itemize}

\subsection{Cierre Técnico}

\begin{itemize}
    \item Pipeline completo:
    \begin{itemize}
        \item Validación cliente $\rightarrow$ JWT $\rightarrow$ validación esquema $\rightarrow$ persistencia.
    \end{itemize}
    \item Propiedades:
    \begin{itemize}
        \item Escalabilidad.
        \item Seguridad.
        \item Bajo acoplamiento.
    \end{itemize}
\end{itemize}


\end{document}